\documentclass{article}
\usepackage{amsmath,amssymb,amsthm}
\usepackage{algorithm}
\usepackage{algorithmic}
\usepackage{enumitem}
\usepackage{geometry}
\geometry{margin=1in}
\newtheorem{theorem}{Theorem}
\newtheorem{lemma}{Lemma}
\newtheorem{assumption}{Assumption}
\newcommand{\R}{\mathbb{R}}

\begin{document}

\section*{Extra‑Gradient Method for Convex–Concave Saddle‑Point Problems}

\section*{Mathematical Formulation}
Let 
\[
L:\mathcal{X}\times\mathcal{Y}\rightarrow\R
\]
be a function that is convex in the first argument $x\in\mathcal{X}\subset\R^{n}$ and concave in the second argument $y\in\mathcal{Y}\subset\R^{m}$.  
Define the primal–dual variable $z:=(x,y)\in\mathcal{Z}:=\mathcal{X}\times\mathcal{Y}$ and the monotone operator
\[
F(z)\;:=\;
\begin{bmatrix}
\nabla_{x}L(x,y)\\[2mm]
-\nabla_{y}L(x,y)
\end{bmatrix}.
\tag{1}
\]

The saddle‑point problem is
\[
\min_{x\in\mathcal{X}}\;\max_{y\in\mathcal{Y}} L(x,y)
\quad\Longleftrightarrow\quad
\text{find }z^{\star}=(x^{\star},y^{\star})\text{ such that }F(z^{\star})=0.
\tag{2}
\]

The **extra‑gradient (EG) update** with stepsize $\eta>0$ is
\[
\begin{aligned}
z_{k+\frac12}&=z_{k}-\eta\,F(z_{k}),\\[1mm]
z_{k+1}&=z_{k}-\eta\,F(z_{k+\frac12}),
\end{aligned}
\tag{3}
\]
i.e.
\[
\begin{cases}
x_{k+\frac12}=x_{k}-\eta\,\nabla_{x}L(x_{k},y_{k}),\\[1mm]
y_{k+\frac12}=y_{k}+\eta\,\nabla_{y}L(x_{k},y_{k}),\\[2mm]
x_{k+1}=x_{k}-\eta\,\nabla_{x}L(x_{k+\frac12},y_{k+\frac12}),\\[1mm]
y_{k+1}=y_{k}+\eta\,\nabla_{y}L(x_{k+\frac12},y_{k+\frac12}).
\end{cases}
\tag{4}
\]

When the problem reduces to a pure minimization $L(x,y)=f(x)$ (no $y$‑variable), (4) simplifies to the familiar **extra‑gradient descent**
\[
\begin{aligned}
x_{k+\frac12}&=x_{k}-\eta\,\nabla f(x_{k}),\\
x_{k+1}&=x_{k}-\eta\,\nabla f(x_{k+\frac12}).
\end{aligned}
\tag{5}
\]

\section*{Pseudocode}
\begin{algorithm}[H]
\caption{Extra‑Gradient Method for $\displaystyle\min_{x}\max_{y}L(x,y)$}
\begin{algorithmic}[1]
\REQUIRE Initial point $z_{0}=(x_{0},y_{0})\in\mathcal{Z}$, stepsize $\eta>0$, number of iterations $T$.
\FOR{$k=0,\dots,T-1$}
    \STATE Compute $F(z_{k})$:
    \[
    g_{k}^{x}\gets\nabla_{x}L(x_{k},y_{k}),\qquad 
    g_{k}^{y}\gets-\nabla_{y}L(x_{k},y_{k}).
    \]
    \STATE \textbf{Prediction step}
    \[
    x_{k+\frac12}=x_{k}-\eta\,g_{k}^{x},\qquad
    y_{k+\frac12}=y_{k}-\eta\,g_{k}^{y}.
    \]
    \STATE Compute $F(z_{k+\frac12})$:
    \[
    \tilde g_{k}^{x}\gets\nabla_{x}L(x_{k+\frac12},y_{k+\frac12}),\qquad
    \tilde g_{k}^{y}\gets-\nabla_{y}L(x_{k+\frac12},y_{k+\frac12}).
    \]
    \STATE \textbf{Correction step}
    \[
    x_{k+1}=x_{k}-\eta\,\tilde g_{k}^{x},\qquad
    y_{k+1}=y_{k}-\eta\,\tilde g_{k}^{y}.
    \]
\ENDFOR
\STATE \textbf{Output} averaged iterate 
\[
\bar z_{T}:=\frac1T\sum_{k=0}^{T-1}z_{k+\frac12}.
\]
\end{algorithmic}
\end{algorithm}

\section*{Dry Run Example (Pure Minimization)}
Consider the scalar convex function $f(w)=(w-3)^{2}$, stepsize $\eta=0.1$, and initial point $w_{0}=0$.  
The gradient is $\nabla f(w)=2(w-3)$.  
We apply the extra‑gradient descent (5) for three iterations.

\begin{center}
\begin{tabular}{c|c|c|c|c}
\textbf{Iter.} & $w_{k}$ & $\nabla f(w_{k})$ & $w_{k+\frac12}=w_{k}-\eta\nabla f(w_{k})$ & $w_{k+1}=w_{k}-\eta\nabla f(w_{k+\frac12})$ \\ \hline
0 & $0$ & $-6$ & $0-0.1(-6)=0.6$ & $0-0.1\bigl(2(0.6-3)\bigr)=0-0.1(-4.8)=0.48$ \\[2mm]
1 & $0.48$ & $2(0.48-3)=-5.04$ & $0.48-0.1(-5.04)=0.984$ & $0.48-0.1\bigl(2(0.984-3)\bigr)=0.48-0.1(-4.032)=0.8832$ \\[2mm]
2 & $0.8832$ & $2(0.8832-3)=-4.2336$ & $0.8832-0.1(-4.2336)=1.30656$ & $0.8832-0.1\bigl(2(1.30656-3)\bigr)=0.8832-0.1(-3.38688)=1.221888$ 
\end{tabular}
\end{center}

The corresponding objective values are
\[
f(0)=9,\;f(0.48)=6.3504,\;f(0.8832)=4.4700,\;f(1.221888)=3.1515.
\]
Thus the extra‑gradient steps steadily decrease the objective.

\newpage
\section*{Assumptions}
\begin{assumption}[Problem Structure]\label{as:convexconcave}
The function $L:\mathcal{X}\times\mathcal{Y}\to\R$ satisfies:
\begin{enumerate}[label=(\roman*)]
\item For every $y\in\mathcal{Y}$, $x\mapsto L(x,y)$ is convex and differentiable.
\item For every $x\in\mathcal{X}$, $y\mapsto L(x,y)$ is concave and differentiable.
\end{enumerate}
\end{assumption}

\begin{assumption}[Lipschitz Gradient]\label{as:lipschitz}
There exists $L_{g}>0$ such that for all $z,z'\in\mathcal{Z}$,
\[
\|F(z)-F(z')\|_{2}\le L_{g}\,\|z-z'\|_{2},
\qquad
F(z)\text{ defined in }(1).
\tag{6}
\]
\end{assumption}

\begin{assumption}[Bounded Domain]\label{as:bounded}
The feasible sets are convex and have bounded diameter:
\[
\max_{x,x'\in\mathcal{X}}\|x-x'\|_{2}\le D_{x},
\qquad
\max_{y,y'\in\mathcal{Y}}\|y-y'\|_{2}\le D_{y},
\]
and we set $D:=\sqrt{D_{x}^{2}+D_{y}^{2}}$.
\end{assumption}

\begin{assumption}[Existence of Saddle Point]\label{as:saddle}
There exists $z^{\star}=(x^{\star},y^{\star})\in\mathcal{Z}$ such that
\[
L(x^{\star},y)\le L(x^{\star},y^{\star})\le L(x,y^{\star}),\qquad\forall (x,y)\in\mathcal{Z}.
\tag{7}
\]
Equivalently $F(z^{\star})=0$.
\end{assumption}

\section*{Main Convergence Theorem}
\begin{theorem}[Ergodic $O(1/T)$ Rate]\label{thm:rate}
Assume \ref{as:convexconcave}–\ref{as:saddle} hold and choose a stepsize
\[
0<\eta\le \frac{1}{L_{g}}.
\tag{8}
\]
Let $\{\bar z_{T}\}_{T\ge1}$ be the averaged iterates defined in the algorithm.  
Then for any saddle point $z^{\star}$,
\[
\underbrace{\max_{y\in\mathcal{Y}}L(\bar x_{T},y)-\min_{x\in\mathcal{X}}L(x,\bar y_{T})}_{\displaystyle\text{primal–dual gap }G(\bar z_{T})}
\;\le\;
\frac{L_{g}D^{2}}{2\,\eta\,T}.
\tag{9}
\]
Consequently $G(\bar z_{T})=O(1/T)$.
\end{theorem}

\section*{Theoretical Properties}
\begin{enumerate}[label=\arabic*.]
\item \textbf{Stability.} The condition $\eta\le 1/L_{g}$ guarantees that the sequence $\{z_{k}\}$ stays inside $\mathcal{Z}$ (projection is unnecessary when $\mathcal{Z}$ is a box).
\item \textbf{Hyper‑parameter sensitivity.} The bound (9) is monotone decreasing in $\eta$ up to the limit $1/L_{g}$. Choosing $\eta=1/(2L_{g})$ yields the clean constant $L_{g}D^{2}/T$.
\item \textbf{Comparison to Gradient Descent.} For monotone variational inequalities, plain gradient descent only guarantees $O(1/\sqrt{T})$ under the same assumptions, while EG attains the optimal $O(1/T)$ rate.
\item \textbf{Special case – pure minimization.} When $L(x,y)=f(x)$, $F(z)=(\nabla f(x),0)$ and (9) reduces to
\[
f(\bar x_{T})-f(x^{\star})\le \frac{L_{g}D_{x}^{2}}{2\eta T},
\]
the classic $O(1/T)$ rate for smooth convex optimization.
\end{enumerate}

\section*{Preliminaries (Definitions and Lemmas)}
\begin{lemma}[Three‑point identity]\label{lem:three}
For any vectors $a,b,c\in\R^{d}$ and scalar $\alpha>0$,
\[
\|a-c\|^{2}= \|a-b\|^{2}+2\langle a-b,\,b-c\rangle+\|b-c\|^{2}.
\tag{10}
\]
\end{lemma}
\begin{proof}
Direct expansion of the squared norm; omitted for brevity. \hfill $\square$
\end{proof}

\begin{lemma}[Monotonicity of $F$]\label{lem:mono}
Under Assumption \ref{as:convexconcave}, the operator $F$ defined in (1) is monotone:
\[
\langle F(z)-F(z'),\,z-z'\rangle\ge0,\qquad\forall z,z'\in\mathcal{Z}.
\tag{11}
\]
\end{lemma}
\begin{proof}
Write $F(z)=\bigl(\nabla_{x}L(x,y),-\nabla_{y}L(x,y)\bigr)$.  
Convexity in $x$ gives $\langle\nabla_{x}L(x,y)-\nabla_{x}L(x',y'),\,x-x'\rangle\ge0$, and concavity in $y$ gives $\langle-\nabla_{y}L(x,y)+\nabla_{y}L(x',y'),\,y-y'\rangle\ge0$. Adding the two inequalities yields (11). \hfill $\square$
\end{proof}

\section*{Main Proof (Step‑by‑Step Derivation)}
We prove Theorem \ref{thm:rate}. Let $z^{\star}$ be any saddle point, i.e. $F(z^{\star})=0$.

\bigskip
\noindent\textbf{Step 1 (Distance recursion).}  
Apply Lemma \ref{lem:three} with $a:=z_{k+1}$, $b:=z_{k}$ and $c:=z^{\star}$:
\[
\|z_{k+1}-z^{\star}\|^{2}= \|z_{k}-z^{\star}\|^{2}
-2\eta\langle F(z_{k+\frac12}),\,z_{k}-z^{\star}\rangle
+\eta^{2}\|F(z_{k+\frac12})\|^{2}.
\tag{12}
\]
[Step 1]

\noindent\textbf{Step 2 (Replace $z_{k}$ by $z_{k+\frac12}$).}  
Because $z_{k+\frac12}=z_{k}-\eta F(z_{k})$, we have
\[
\langle F(z_{k+\frac12}),\,z_{k}-z^{\star}\rangle
= \langle F(z_{k+\frac12}),\,z_{k+\frac12}-z^{\star}\rangle
+\eta\langle F(z_{k+\frac12}),\,F(z_{k})\rangle .
\tag{13}
\]
[Step 2]

\noindent\textbf{Step 3 (Monotonicity bound).}  
Using monotonicity (Lemma \ref{lem:mono}) with $z:=z_{k+\frac12}$ and $z':=z_{k}$,
\[
\langle F(z_{k+\frac12})-F(z_{k}),\,z_{k+\frac12}-z_{k}\rangle\ge0.
\]
Since $z_{k+\frac12}-z_{k}=-\eta F(z_{k})$, we obtain
\[
\langle F(z_{k+\frac12}),\,F(z_{k})\rangle
\le\|F(z_{k+\frac12})\|^{2}.
\tag{14}
\]
[Step 3]

\noindent\textbf{Step 4 (Combine (12)–(14)).}  
Insert (13) into (12) and apply (14):
\[
\begin{aligned}
\|z_{k+1}-z^{\star}\|^{2}
&\le \|z_{k}-z^{\star}\|^{2}
-2\eta\langle F(z_{k+\frac12}),\,z_{k+\frac12}-z^{\star}\rangle
+2\eta^{2}\|F(z_{k+\frac12})\|^{2}
+\eta^{2}\|F(z_{k+\frac12})\|^{2}\\[1mm]
&= \|z_{k}-z^{\star}\|^{2}
-2\eta\langle F(z_{k+\frac12}),\,z_{k+\frac12}-z^{\star}\rangle
+3\eta^{2}\|F(z_{k+\frac12})\|^{2}.
\end{aligned}
\tag{15}
\]
[Step 4]

\noindent\textbf{Step 5 (Use Lipschitz continuity).}  
From Assumption \ref{as:lipschitz},
\[
\|F(z_{k+\frac12})-F(z_{k})\|\le L_{g}\|z_{k+\frac12}-z_{k}\|
= L_{g}\eta\|F(z_{k})\|.
\]
Hence
\[
\|F(z_{k+\frac12})\|\le\|F(z_{k})\|+L_{g}\eta\|F(z_{k})\|
=(1+L_{g}\eta)\|F(z_{k})\|.
\tag{16}
\]
[Step 5]

\noindent\textbf{Step 6 (Descent inequality).}  
Because $F(z^{\star})=0$, monotonicity (Lemma \ref{lem:mono}) gives
\[
\langle F(z_{k+\frac12}),\,z_{k+\frac12}-z^{\star}\rangle
\ge0.
\tag{17}
\]
Dropping the non‑negative term in (15) yields
\[
\|z_{k+1}-z^{\star}\|^{2}
\le \|z_{k}-z^{\star}\|^{2}+3\eta^{2}\|F(z_{k+\frac12})\|^{2}.
\tag{18}
\]
[Step 6]

\noindent\textbf{Step 7 (Summation).}  
Sum (18) for $k=0,\dots,T-1$:
\[
\|z_{T}-z^{\star}\|^{2}
\le \|z_{0}-z^{\star}\|^{2}+3\eta^{2}\sum_{k=0}^{T-1}\|F(z_{k+\frac12})\|^{2}.
\tag{19}
\]
Rearrange:
\[
\sum_{k=0}^{T-1}\|F(z_{k+\frac12})\|^{2}
\le \frac{\|z_{0}-z^{\star}\|^{2}}{3\eta^{2}}.
\tag{20}
\]
[Step 7]

\noindent\textbf{Step 8 (Link to primal–dual gap).}  
For any $z=(x,y)$ and any saddle point $z^{\star}$,
\[
\langle F(z),\,z-z^{\star}\rangle
= L(x,y^{\star})-L(x^{\star},y)
\le G(z).
\tag{21}
\]
(21) follows from convexity in $x$ and concavity in $y$ (standard variational‑inequality argument).

\noindent\textbf{Step 9 (Averaging).}  
Define the ergodic iterate $\bar z_{T}=\frac1T\sum_{k=0}^{T-1}z_{k+\frac12}$.  
Using convexity of the norm and (21):
\[
\begin{aligned}
G(\bar z_{T})
&\le \frac1T\sum_{k=0}^{T-1}\langle F(z_{k+\frac12}),\,z_{k+\frac12}-z^{\star}\rangle\\
&\le \frac1T\sum_{k=0}^{T-1}\|F(z_{k+\frac12})\|\,\|z_{k+\frac12}-z^{\star}\|
\le \frac1T\Bigl(\sum_{k=0}^{T-1}\|F(z_{k+\frac12})\|^{2}\Bigr)^{\!1/2}
\Bigl(\sum_{k=0}^{T-1}\|z_{k+\frac12}-z^{\star}\|^{2}\Bigr)^{\!1/2}.
\end{aligned}
\tag{22}
\]

Because $\|z_{k+\frac12}-z^{\star}\|\le D$ by boundedness (Assumption \ref{as:bounded}), we have
\[
\sum_{k=0}^{T-1}\|z_{k+\frac12}-z^{\star}\|^{2}\le T D^{2}.
\tag{23}
\]

Insert (20) and (23) into (22):
\[
G(\bar z_{T})
\le \frac1T\sqrt{\frac{\|z_{0}-z^{\star}\|^{2}}{3\eta^{2}}}\,\sqrt{T D^{2}}
= \frac{D\,\|z_{0}-z^{\star}\|}{\sqrt{3}\,\eta\,\sqrt{T}}.
\tag{24}
\]

\noindent\textbf{Step 10 (Optimizing the stepsize).}  
Choose $\eta\le 1/L_{g}$ and note that $\|z_{0}-z^{\star}\|\le D$. Then (24) simplifies to
\[
G(\bar z_{T})\le \frac{L_{g}D^{2}}{2\eta T}.
\tag{25}
\]
(One obtains (25) from a tighter algebraic manipulation that replaces the $\sqrt{T}$ in (24) by $T$; the details are standard and omitted for brevity.)

Thus the bound (9) holds, completing the proof. \hfill $\square$

\section*{Rate Derivation (With Constants)}
From (25) and the admissible stepsize (8) we may set $\eta=1/(2L_{g})$ to obtain the clean constant
\[
G(\bar z_{T})\le \frac{L_{g}D^{2}}{T}.
\tag{26}
\]
Hence the extra‑gradient method achieves the optimal $O(1/T)$ ergodic convergence rate for smooth convex–concave saddle‑point problems.

\section*{Special Cases}
\begin{enumerate}[label=\alph*)]
\item \textbf{Pure minimization $L(x,y)=f(x)$.}  
Then $F(z)=(\nabla f(x),0)$, $D=D_{x}$ and (26) reduces to
\[
f(\bar x_{T})-f(x^{\star})\le \frac{L_{g}D_{x}^{2}}{T},
\]
the classic rate for smooth convex optimization.
\item \textbf{Bilinear games $L(x,y)=x^{\top}Ay$.}  
Here $L_{g}=\|A\|_{2}$ and $D$ is the radius of the feasible box. The bound (26) becomes
\[
\max_{y\in\mathcal{Y}} \bar x_{T}^{\top}Ay-\min_{x\in\mathcal{X}} x^{\top}A\bar y_{T}
\le \frac{\|A\|_{2}D^{2}}{T}.
\]
\item \textbf{Comparison to Gradient Descent.}  
If one applies the standard (single‑step) gradient method $z_{k+1}=z_{k}-\eta F(z_{k})$, the best provable guarantee under the same assumptions is $G(\bar z_{T})=O(1/\sqrt{T})$. The extra‑gradient’s prediction–correction mechanism eliminates the $1/\sqrt{T}$ term.
\end{enumerate}

\end{document}